% GENERATED FILE. Edit canonical lessons, then run npm run generate:books.
% canonical-source-hash: fnv1a64:32d3dd635f579974
% canonical-lessons: KA-C431-tirmanisu, KA-C431-samalocisu, KA-C431-samvada, KA-C431-tharavu, KA-C431-olagollu

\chapter{Reaching a Decision}
\label{ch:ka-reaching-a-decision}
\hlchaptermodality{431}

\begin{chapteropening}
\textbf{By the end of this chapter:} \emph{I can say to conclude, to confer, a dialogue, a formal resolution and to include.}

Complete the last lesson of chapter 431: I can say to conclude, to confer, a dialogue, a formal resolution and to include.
\end{chapteropening}

\section[\texorpdfstring{\kn{ತೀರ್ಮಾನಿಸು} (tīrmānisu)}{ತೀರ್ಮಾನಿಸು (tīrmānisu)}]{\kn{ತೀರ್ಮಾನಿಸು} (tīrmānisu) — to conclude; to resolve}

\label{lesson:KA-C431-tirmanisu}



\begin{quote}
Before the new one: say the Kannada for a stipulation, then the Kannada for a requirement.
\end{quote}

\subsection*{You'll want to know: \kn{ತೀರ್ಮಾನಿಸು}}
\textbf{\kn{ತೀರ್ಮಾನಿಸು}} — \emph{tīrmānisu} — \textquotedblleft{}to conclude; to resolve\textquotedblright{}.

\begin{quote}
\kn{ಕೊನೆಗೆ} \kn{ನಾನು} \kn{ಊರಿಗೆ} \kn{ಹಿಂತಿರುಗಲು} \kn{ತೀರ್ಮಾನಿಸಿದೆ}.
\end{quote}

\begin{quote}
\emph{konege nānu ūrige hintirugalu tīrmāniside.}
\end{quote}

\begin{quote}
In the end I resolved to go back to my town.
\end{quote}

Past (I …): \kn{ತೀರ್ಮಾನಿಸಿದೆ}, \emph{tīrmāniside}.

\kn{ತೀರ್ಮಾನ}, a decision or a conclusion.

\begin{cousinweb}
\textbf{\kn{ತೀರ್ಮಾನಿಸು}} (\emph{tīrmānisu}) is built on \emph{tīrmāna}, \textquotedblleft{}a settlement, a conclusion\textquotedblright{}, from native \emph{tīr}, \textquotedblleft{}to be finished, to be settled\textquotedblright{} (DEDR 3278).
\end{cousinweb}

\begin{grammarlens}[title={What you've built}]
One more word for saying what happens if.
\end{grammarlens}

\subsection*{Your turn}
\begin{itemize}
\raggedright
  \item \emph{Say it:} \emph{tīrmānisu}
  \item \emph{Say it:} \emph{tīrmānisu}, once more
  \item \emph{Say it:} \emph{konege nānu ūrige hintirugalu tīrmāniside.} — In the end I resolved to go back to my town.
  \item \emph{Recall:} say the Kannada for a stipulation, then the Kannada for a requirement, then say \emph{tīrmānisu} again
\end{itemize}

\begin{tcolorbox}[breakable,colback=teal!4,colframe=teal!35!black,title={Before you move on}]
What does \kn{ತೀರ್ಮಾನಿಸು} mean? (\textquotedblleft{}To conclude; to resolve\textquotedblright{}.) Say it once more.
\end{tcolorbox}
\section[\texorpdfstring{\kn{ಸಮಾಲೋಚಿಸು} (samālōcisu)}{ಸಮಾಲೋಚಿಸು (samālōcisu)}]{\kn{ಸಮಾಲೋಚಿಸು} (samālōcisu) — to confer; to consult}

\label{lesson:KA-C431-samalocisu}



\begin{quote}
Before the new one: say the Kannada for a requirement, then the Kannada for to conclude.
\end{quote}

\subsection*{You'll want to know: \kn{ಸಮಾಲೋಚಿಸು}}
\textbf{\kn{ಸಮಾಲೋಚಿಸು}} — \emph{samālōcisu} — \textquotedblleft{}to confer; to consult\textquotedblright{}.

\begin{quote}
\kn{ನಿರ್ಧಾರ} \kn{ತೆಗೆದುಕೊಳ್ಳುವ} \kn{ಮೊದಲು} \kn{ನಾನು} \kn{ಅಪ್ಪನ} \kn{ಜೊತೆ} \kn{ಸಮಾಲೋಚಿಸಿದೆ}.
\end{quote}

\begin{quote}
\emph{nirdhāra tegedukoḷḷuva modalu nānu appana jote samālōciside.}
\end{quote}

\begin{quote}
Before deciding I conferred with Father.
\end{quote}

Past (I …): \kn{ಸಮಾಲೋಚಿಸಿದೆ}, \emph{samālōciside}.

With whom: -\kn{ಜೊತೆ}. \kn{ಸಮಾಲೋಚನೆ}, a consultation.

\begin{cousinweb}
\textbf{\kn{ಸಮಾಲೋಚಿಸು}} (\emph{samālōcisu}) is built on Sanskrit \emph{sam-}, \textquotedblleft{}together\textquotedblright{}, and \emph{ālocana}, \textquotedblleft{}seeing, considering\textquotedblright{} (Monier-Williams), with the Kannada verb ending -\emph{isu}.
\end{cousinweb}

\begin{grammarlens}[title={What you've built}]
One more word for saying what happens if.
\end{grammarlens}

\subsection*{Your turn}
\begin{itemize}
\raggedright
  \item \emph{Say it:} \emph{samālōcisu}
  \item \emph{Say it:} \emph{samālōcisu}, once more
  \item \emph{Say it:} \emph{nirdhāra tegedukoḷḷuva modalu nānu appana jote samālōciside.} — Before deciding I conferred with Father.
  \item \emph{Recall:} say the Kannada for a requirement, then the Kannada for to conclude, then say \emph{samālōcisu} again
\end{itemize}

\begin{tcolorbox}[breakable,colback=teal!4,colframe=teal!35!black,title={Before you move on}]
What does \kn{ಸಮಾಲೋಚಿಸು} mean? (\textquotedblleft{}To confer; to consult\textquotedblright{}.) Say it once more.
\end{tcolorbox}
\section[\texorpdfstring{\kn{ಸಂವಾದ} (saṃvāda)}{ಸಂವಾದ (saṃvāda)}]{\kn{ಸಂವಾದ} (saṃvāda) — a dialogue; a discussion}

\label{lesson:KA-C431-samvada}



\begin{quote}
Before the new one: say the Kannada for to conclude, then the Kannada for to confer.
\end{quote}

\subsection*{You'll want to know: \kn{ಸಂವಾದ}}
\textbf{\kn{ಸಂವಾದ}} — \emph{saṃvāda} — \textquotedblleft{}a dialogue; a discussion\textquotedblright{}.

\begin{quote}
\kn{ವಿದ್ಯಾರ್ಥಿಗಳು} \kn{ಲೇಖಕರ} \kn{ಜೊತೆ} \kn{ಸಂವಾದ} \kn{ನಡೆಸಿದರು}.
\end{quote}

\begin{quote}
\emph{vidyārthigaḷu lēkhakara jote saṃvāda naḍesidaru.}
\end{quote}

\begin{quote}
The students held a discussion with the author.
\end{quote}

\kn{ಸಂವಾದ} \kn{ನಡೆಸು}, to hold a dialogue.

\begin{cousinweb}
\textbf{\kn{ಸಂವಾದ}} (\emph{saṃvāda}) is borrowed from Sanskrit \emph{saṃvāda}, \textquotedblleft{}speaking together, conversation\textquotedblright{} (Monier-Williams).
\end{cousinweb}

\begin{grammarlens}[title={What you've built}]
One more word for saying what happens if.
\end{grammarlens}

\subsection*{Your turn}
\begin{itemize}
\raggedright
  \item \emph{Say it:} \emph{saṃvāda}
  \item \emph{Say it:} \emph{saṃvāda}, once more
  \item \emph{Say it:} \emph{vidyārthigaḷu lēkhakara jote saṃvāda naḍesidaru.} — The students held a discussion with the author.
  \item \emph{Recall:} say the Kannada for to conclude, then the Kannada for to confer, then say \emph{saṃvāda} again
\end{itemize}

\begin{tcolorbox}[breakable,colback=teal!4,colframe=teal!35!black,title={Before you move on}]
What does \kn{ಸಂವಾದ} mean? (\textquotedblleft{}A dialogue; a discussion\textquotedblright{}.) Say it once more.
\end{tcolorbox}
\section[\texorpdfstring{\kn{ಠರಾವು} (ṭharāvu)}{ಠರಾವು (ṭharāvu)}]{\kn{ಠರಾವು} (ṭharāvu) — a formal resolution}

\label{lesson:KA-C431-tharavu}



\begin{quote}
Before the new one: say the Kannada for to confer, then the Kannada for a dialogue.
\end{quote}

\subsection*{You'll want to know: \kn{ಠರಾವು}}
\textbf{\kn{ಠರಾವು}} — \emph{ṭharāvu} — \textquotedblleft{}a formal resolution\textquotedblright{}.

\begin{quote}
\kn{ಸಭೆ} \kn{ಹೊಸ} \kn{ಶಾಲೆಯ} \kn{ಪರವಾಗಿ} \kn{ಠರಾವು} \kn{ಮಾಡಿತು}.
\end{quote}

\begin{quote}
\emph{sabhe hosa śāleya paravāgi ṭharāvu māḍitu.}
\end{quote}

\begin{quote}
The meeting passed a resolution in favour of a new school.
\end{quote}

\kn{ಠರಾವು} \kn{ಮಾಡು} / \kn{ಅಂಗೀಕರಿಸು}, to pass a resolution.

\begin{cousinweb}
\textbf{\kn{ಠರಾವು}} (\emph{ṭharāvu}) is usually traced to Hindustani or Marathi \emph{ṭharāv}, \textquotedblleft{}a fixing, a resolution\textquotedblright{}; DEDR and CDIAL do not list the Kannada form.
\end{cousinweb}

\begin{grammarlens}[title={What you've built}]
One more word for saying what happens if.
\end{grammarlens}

\subsection*{Your turn}
\begin{itemize}
\raggedright
  \item \emph{Say it:} \emph{ṭharāvu}
  \item \emph{Say it:} \emph{ṭharāvu}, once more
  \item \emph{Say it:} \emph{sabhe hosa śāleya paravāgi ṭharāvu māḍitu.} — The meeting passed a resolution in favour of a new school.
  \item \emph{Recall:} say the Kannada for to confer, then the Kannada for a dialogue, then say \emph{ṭharāvu} again
\end{itemize}

\begin{tcolorbox}[breakable,colback=teal!4,colframe=teal!35!black,title={Before you move on}]
What does \kn{ಠರಾವು} mean? (\textquotedblleft{}A formal resolution\textquotedblright{}.) Say it once more.
\end{tcolorbox}
\section[\texorpdfstring{\kn{ಒಳಗೊಳ್ಳು} (oḷagoḷḷu)}{ಒಳಗೊಳ್ಳು (oḷagoḷḷu)}]{\kn{ಒಳಗೊಳ್ಳು} (oḷagoḷḷu) — to include}

\label{lesson:KA-C431-olagollu}



\begin{quote}
Before the new one: say the Kannada for a dialogue, then the Kannada for a formal resolution.
\end{quote}

\subsection*{You'll want to know: \kn{ಒಳಗೊಳ್ಳು}}
\textbf{\kn{ಒಳಗೊಳ್ಳು}} — \emph{oḷagoḷḷu} — \textquotedblleft{}to include\textquotedblright{}.

\begin{quote}
\kn{ಪ್ರವಾಸದ} \kn{ಬೆಲೆ} \kn{ಊಟವನ್ನೂ} \kn{ಒಳಗೊಂಡಿತ್ತು}.
\end{quote}

\begin{quote}
\emph{pravāsada bele ūṭavannū oḷagoṇḍittu.}
\end{quote}

\begin{quote}
The price of the trip included meals.
\end{quote}

Past (it …): \kn{ಒಳಗೊಂಡಿತು}, \emph{oḷagoṇḍitu}.

Mostly of things: \kn{ಒಳಗೊಂಡಿದೆ}, it includes.

\begin{cousinweb}
\textbf{\kn{ಒಳಗೊಳ್ಳು}} (\emph{oḷagoḷḷu}) is native Dravidian, built from \emph{oḷ}, \textquotedblleft{}the inside\textquotedblright{} (DEDR 698), and \emph{koḷ}, \textquotedblleft{}to take\textquotedblright{} (DEDR 2151): to take inside.
\end{cousinweb}

\begin{grammarlens}[title={What you've built}]
One more word for saying what happens if.
\end{grammarlens}

\subsection*{Your turn}
\begin{itemize}
\raggedright
  \item \emph{Say it:} \emph{oḷagoḷḷu}
  \item \emph{Say it:} \emph{oḷagoḷḷu}, once more
  \item \emph{Say it:} \emph{pravāsada bele ūṭavannū oḷagoṇḍittu.} — The price of the trip included meals.
  \item \emph{Recall:} say the Kannada for a dialogue, then the Kannada for a formal resolution, then say \emph{oḷagoḷḷu} again
\end{itemize}

\begin{tcolorbox}[breakable,colback=teal!4,colframe=teal!35!black,title={Before you move on}]
What does \kn{ಒಳಗೊಳ್ಳು} mean? (\textquotedblleft{}To include\textquotedblright{}.) Say it once more.
\end{tcolorbox}
